El crecimiento exponencial de los datos astronómicos ha planteado nuevos desafíos en el análisis eficiente de curvas de luz, fundamentales para el estudio de estrellas variables periódicas. Este trabajo aborda la optimización y modernización de la biblioteca \acrfull{feets}, una herramienta de Python para la extracción de características de curvas de luz. Presentamos una refactorización integral que incorpora una nueva implementación paralela que busca reducir los tiempos de cómputo, con un rediseño arquitectónico basado en principios de ingeniería de \emph{software}. Los resultados obtenidos demuestran una reducción notable en los tiempos de ejecución y una mejora sustancial en la calidad del código, lo que posiciona a la nueva versión de \acrshort{feets} como una herramienta escalable para el análisis astronómico de grandes volúmenes de datos.